\section{Introduction}\label{sec:tsps-introduction}

\paragraph{Certifiers as participants.} Transparency logs, permissioned ledgers
and privacy preserving payment systems each contain an entity that certifies
state on behalf of everyone else: a log provider, an endorser, a ledger
administrator. These entities are not artefacts of poor design that a better
protocol would remove. They exist because something has to establish that a
certificate was issued, that a transaction was valid, or that an account held a
particular balance at a particular time. The question this chapter asks is how
narrowly the capability each of them holds can be defined.

Thresholdising the certifier divides the authority to certify, so that no single
participant can produce a certification alone and no single key compromise is
enough to forge one. Structure preservation constrains something different, the
knowledge that accompanies that authority. Because signing acts directly on
group elements, a certifier can sign a commitment it is unable to open, and the
party holding the resulting signature later proves statements about the
committed value without revealing it. Authority and knowledge are separated, and
the institution performs its operational function without being trusted with the
information it certifies. What remains is practicality, by which we mean
minimising the interaction, coordination, communication and computational costs
that the protocol imposes on its participants. It is for this reason that our
constructions are designed so that the online phase of signing requires no
interaction between the certifiers at all.

\paragraph{Structure preserving signatures.} Structure preserving signatures are
digital signatures in which the messages, public keys, and signatures are all
group elements, and verification consists solely of checking pairing product
equations~\cite{C:AFGHO10,C:AGHO11,TCC:AGOT14,C:KilPanWee15,AC:Groth15,RSA:Ghadafi16}.
Coupled with non-interactive zero knowledge proofs, structure preserving
signatures have proved useful across a range of privacy preserving applications,
including blind signatures, delegatable anonymous credentials, and group
signatures, to name a
few~\cite{AFRICACRYPT:FucVer10,CCS:CamDriDub17,RSA:CriLys19,PoPETS:MSBM23,AC:Groth06,AC:Groth07,C:LibPetYun15}.

Structure preserving signatures also find use in privacy preserving payment
systems built over a shared ledger. Currency is represented by a token encoding
the owner, value, and a unique serial
number~\cite{EPRINT:ACDDET19,EPRINT:TBAAGP22,CCS:KiaKohSar22}. For privacy, the
shared ledger records a hiding commitment to the token, generally a Pedersen
commitment, rather than the token itself~\cite{C:Pedersen91}. To transfer
ownership privately, the owner interacts with a certifier to obtain a blind
Pointcheval-Sanders signature on the commitment's opening. That signature is
later used to prove that the token being spent is recorded on the ledger, and,
via its serial number, that it has not been spent
before~\cite{RSA:PoiSan16,RSA:PoiSan18,NDSS:SABMD19}. Structure preserving
signatures remove the need for this interaction: since signing acts directly on
the group elements that make up the commitment, the certifier can sign tokens as
they appear on the ledger and simply publish the resulting signatures.  Removing
the interaction is a gain in more than efficiency. A blind issuance protocol
gives the certifier a per-transfer channel to each spender, which can be timed,
logged and correlated even when the value being signed stays hidden.  Publishing
signatures over ledger state removes that channel altogether.

A similar idea applies where the ledger stores commitments to user accounts,
pairs (owner, holdings), rather than committed tokens. Such accounts form the
basis for periodic reporting on user holdings, for instance proving that
holdings do not exceed some threshold when filing taxes. Since tax authorities
typically lack direct access to the shared ledger, they instead rely on
signatures produced by the ledger's administrators. Structure preserving
signatures let the administrators sign the committed accounts together with a
timestamp and publish the result directly, without involving the user. The user
can then use the published signature to prove properties of their holdings
without disclosing anything beyond what is required.

\paragraph{Transparency logs.} Certificate transparency (CT) and key
transparency (KT) are the main examples of transparency logs. In both cases, a
transparency log is an append-only, publicly verifiable record, of either
certificates or keys. Each entry is sent to a log provider, signed, and appended
to the log. In CT systems, the log server returns a Signed Certificate Timestamp
(SCT) to the certificate provider, and in KT systems, the mapping from user
identity to public key is what the log server signs.  Transparency logs let
users and auditors verify that the certificates or keys they interact with carry
legitimacy beyond the app serving them, moving trust onto the log providers
instead. Threshold signatures narrow what any one log provider can do,
distributing the capability across several servers. A threshold structure
preserving signature given to the key or certificate owner would let the
structure preserving property of the signature be used to prove, in zero
knowledge, that a key is recorded in the log, for instance for privacy
preserving identity verification. Since the structure preserving property lets
the public key itself be signed directly, without first hashing it, the
resulting zero knowledge proofs are more efficient than they would otherwise be.

\paragraph{Permissioned ledgers.} A permissioned ledger is a distributed ledger
where token issuance, transaction verification, and participation are controlled
by one or a small number of authorities. Permissioned ledgers are of particular
interest for platforms such as Central Bank Digital Currencies (CBDCs), as they
offer a transparent system enabling regulatory oversight while still supporting
transaction privacy and compliance with the monetary policy of the central bank
in question. One way to build such a system has every transaction executed and,
if valid, signed by an endorser before being added to the ledger. The signed
state committed to the ledger is a commitment to the value and type of asset.
Threshold structure preserving signatures would let trust be distributed between
multiple endorsers via the threshold property, while the structure preserving
property lets users perform efficient proofs of statements about their assets on
the ledger, for example for compliance purposes.

In each of these applications the certifier, endorser, log provider or ledger
administrator is not an accident of the design. Something has to attest that a
token is recorded, that a key is in the log, or that a transaction was
validated, and that role is one the system genuinely requires. The question is
how narrowly to define the capability attached to it. A threshold signature
answers that directly. The role stays, but the ability to produce an attestation
is separated from the ability to hold the signing key, so that no single
certifier can act alone and none of them ever holds the key. Structure
preserving signatures are fit for these applications only once they admit an
efficient threshold counterpart. We give two threshold structure preserving
signature schemes, building on two existing structure preserving signatures
together with secure multiparty computation (MPC)
techniques~\cite{C:AGHO11,C:KilPanWee15,C:DPSZ12,CCS:KelOrsSch16,C:GLOPS21}.  We
select the scheme of Abe, Groth, Haralambiev, and Ohkubo for its optimal
signature size, and that of Kiltz, Pan, and Wee for its reliance on standard
assumptions alone~\cite{C:AGHO11,C:KilPanWee15}. Both underlying signatures let
the threshold signers perform the majority of their computation offline, in a
pre-processing phase before the message to be signed is known. This suits
high-throughput applications, such as payments, that experience interspersed
periods of low and high traffic. Since our threshold signatures require
distributed computation over group elements rather than only field elements, we
build on techniques that extend MPC protocols over field elements to protocols
over group elements~\cite{NDSS:GPVRNC23}. Our schemes inherit the security
properties of the underlying structure preserving signatures, and the cost of
the online phase of threshold signing is comparable to that of ordinary signing.

This chapter treats unforgeability rather than privacy directly. Privacy in the
applications above arises from composition: the certifier signs committed state,
and the holder of the signature discloses what it chooses through a zero
knowledge proof. The contribution of a threshold structure preserving signature
to that story is twofold. It divides the authority to certify without giving any
certifier access to the certified value, and, by signing ledger state rather
than blinded requests, it removes the per user interaction that issuance would
otherwise require.

\subsection{Related work}\label{sec:tsps-related}

\paragraph{Structure preserving signatures.} Structure preserving signatures
have messages, signatures, and verification keys that are all group elements.
They were first introduced, and have since been studied extensively, with
matching impossibility results known for signature size in the generic group
model (GGM), from $q$-type assumptions, and from non-$q$-type
assumptions~\cite{C:AFGHO10,C:AGHO11,AC:AbeGroOhk11}.

The original structure preserving signature scheme signs messages of arbitrary
length with signatures of seven group elements, and is secure in the standard
model under a non-interactive assumption~\cite{C:AFGHO10}. Abe, Groth,
Haralambiev, and Ohkubo then established a lower bound of three group elements
for signature size in type III pairing groups, those without an efficiently
computable isomorphism between $\group{1}$ and $\group{2}$, and gave a matching
construction secure in the GGM~\cite{C:AGHO11}. These signatures achieve strong
unforgeability, and remain three group elements regardless of message length. A
linearly homomorphic structure preserving signature, also of just three group
elements, enables further applications~\cite{C:LPJY13}.

Fully structure preserving signatures additionally have secret keys that are
group elements~\cite{JC:AGKOT19}. Structure preserving signatures integrate well
with Groth-Sahai proofs of knowledge, so a group element secret key lets a user
also prove statements about their own key~\cite{EC:GroSah08}. The cost is that
signature length grows with the square root of the message
length~\cite{JC:AGKOT19}.

\begin{table}[t] \centering \resizebox{\textwidth}{!}{ \begin{tabular}{ c | c |
c | c } \toprule & Assumption/model & $|\SPSsig|$ & Notable property \\ \midrule
\cite{C:AFGHO10} & standard model, non-interactive assumption & 7 group elements
& signs messages of arbitrary length \\ \cite{C:AGHO11} & GGM & 3 group elements
& strong unforgeability, size-optimal in type III groups \\ \cite{C:LPJY13} & --
& 3 group elements & linearly homomorphic \\ \cite{JC:AGKOT19} & -- & grows with
$\sqrt{\cdot}$ of message length & fully structure preserving (group element
secret key) \\ \cite{C:KilPanWee15} & DDH & 7 group elements & security from a
standard, non-$q$-type assumption \\ \bottomrule \end{tabular} }
\caption{Structure preserving signature schemes.} \label{tab:tsps-related-sps}
\end{table}

\paragraph{Threshold structure preserving signatures.} Threshold structure
preserving signatures have only recently been considered. One non-interactive
construction is restricted to signing indexed Diffie-Hellman tuples of the form
$(\gen^\msg, \ggen^\msg)$~\cite{AC:CKPSS23}. This restricts the scheme to
messages whose discrete logarithm relative to $\gen$ or $\ggen$ is known. It
therefore rules out signing an existing public key or commitment directly.
either the threshold signers would learn someone else's secret key, or the
element simply could not be signed under the scheme. Another non-interactive
threshold structure preserving signature instead is capable of signing arbitrary
group elements, and is secure under the SXDH assumption~\cite{PKC:MMSST24}. It
builds on the structure preserving signature of Kiltz, Pan, and Wee, and its
signatures consist of seven group elements~\cite{PKC:MMSST24,C:KilPanWee15}.

\begin{table}[t] \centering \begin{tabular}{ c | c | c | c } \toprule &
Assumption & Message space & Built from \\ \midrule \cite{AC:CKPSS23} & PS + ROM
& indexed Diffie-Hellman tuples $(\gen^\msg, \ggen^\msg)$ & -- \\
\cite{PKC:MMSST24} & SXDH & arbitrary $\vec{\msg} \in \group{1}^\ell$ &
\cite{C:KilPanWee15} \\ \Cref{sec:tsps-abe} (ours) & GGM & arbitrary $\vec{\msg}
\in \group{1}^\ell$ & \cite{C:AGHO11} \\ \Cref{sec:tsps-kiltz} (ours) & SXDH &
arbitrary $\vec{\msg} \in \group{1}^\ell$ & \cite{C:KilPanWee15} \\ \bottomrule
\end{tabular} \caption{Threshold structure preserving signature schemes.
Concrete key, signature, and communication sizes for these schemes are compared
in \Cref{tab:tsps-comp} and \Cref{tab:tsps-comms} in
\Cref{sec:tsps-evaluation}.} \label{tab:tsps-related-tsps}
%
\end{table}

\paragraph{Our contribution.} We design two threshold structure preserving
signature schemes relative to an abstract secure multiparty computation
protocol. We specify its interfaces, and require its security as a set of stated
properties rather than as a definitional target our schemes are proved to
realise. These interfaces extend the standard ones over field elements to
secret-shared group elements, which our schemes require and which existing MPC
frameworks do not provide. Computing on shared group elements is itself well
studied, and we claim no novelty for it. \Cref{sec:tsps-mpc} sets out what we
take from existing work and what we
add~\cite{FOCS:Feldman87,IMA:SmaTal19,LC:ADEO21,NDSS:GPVRNC23}. Both schemes
are designed so that the online phase, run once the messages to be signed are
known, requires no interaction between parties, keeping its cost comparable to
that of ordinary, non-threshold signing. We prove the security of the threshold
signatures relative to these properties, and show that they offer the same
security guarantees as their non-threshold counterparts.  Finally, we benchmark
them using two MPC frameworks that realise the required interfaces: ATLAS, which
assumes an honest majority, and MASCOT, which assumes a dishonest
majority~\cite{C:GLOPS21,CCS:KelOrsSch16}. Supporting them required us to extend
MP-SPDZ with the group element interfaces our schemes need~\cite{CCS:Keller20}.
The former suits settings where robustness and high availability are desirable,
the latter settings where resilience to key compromise takes precedence.

We also implement previously suggested threshold structure preserving signature
schemes, so as to compare concrete efficiency. Our scheme built from the
signature of Abe, Groth, Haralambiev, and Ohkubo yields optimally small
threshold signatures, consisting of just three group elements~\cite{C:AGHO11}.

\paragraph{Layout.} \Cref{sec:tsps-definitions} formalises the security
properties of threshold structure preserving signatures, and introduces the MPC
interfaces that our constructions build on. \Cref{sec:tsps-construction} gives a
construction that thresholdises a generic structure preserving signature scheme
using those interfaces. \Cref{sec:tsps-instantiation} instantiates it twice,
with the signature schemes of Abe, Groth, Haralambiev, and Ohkubo and of Kiltz,
Pan, and Wee, extending the MPC interfaces to secret-shared group elements along
the way~\cite{C:AGHO11,C:KilPanWee15}.  Both the construction and its
instantiations are proven secure in \Cref{sec:tsps-security}.
\Cref{sec:tsps-evaluation} evaluates our approach, and \Cref{sec:tsps-summary}
concludes.
